% =====================================================
% 题型：解三角形和差化积选择题 (q_triangle_sum_to_product.tex)
% =====================================================
% =====================================================
% 难度：★★ (中档)
% =====================================================
\newcommand{\qTriangleSumToProductStem}{在 \(\triangle ABC\) 中，若 \(A-B=\dfrac{\pi}{3}\)，且 \(\sin A+\sin B=1\)，则 \(\cos C=\) \hfill （\quad）}

\newcommand{\qTriangleSumToProductOptions}{%
  \mcoptions[4]{\(-\dfrac13\)}{\(\dfrac13\)}{\(-\dfrac79\)}{\(\dfrac79\)}%
}

\newcommand{\qTriangleSumToProductAnswer}{A}

\newcommand{\qTriangleSumToProductSolution}{%
  在 \(\triangle ABC\) 中，已知 \(A-B = \dfrac{\pi}{3}\)，\(\sin A + \sin B = 1\)。
  
  由三角函数和差化积公式：
  \[
    \sin A + \sin B = 2 \sin\frac{A+B}{2} \cos\frac{A-B}{2}.
  \]
  
  代入 \(A-B = \dfrac{\pi}{3}\) 可得：
  \[
    \frac{A-B}{2} = \frac{\pi}{6}.
  \]
  
  因为 \(\cos\dfrac{\pi}{6} = \dfrac{\sqrt{3}}{2}\)，所以：
  \[
    \sin A + \sin B = 2 \sin\frac{A+B}{2} \cdot \frac{\sqrt{3}}{2} = \sqrt{3} \sin\frac{A+B}{2}.
  \]
  
  结合已知等式 \(\sin A + \sin B = 1\)，得：
  \[
    \sqrt{3} \sin\frac{A+B}{2} = 1 \implies \sin\frac{A+B}{2} = \frac{1}{\sqrt{3}}.
  \]
  
  设 \(S = A+B\)，则由二倍角公式可求 \(\cos S\)：
  \[
    \cos S = \cos(A+B) = 1 - 2\sin^2\frac{A+B}{2} = 1 - 2 \cdot \left(\frac{1}{\sqrt{3}}\right)^2 = 1 - \frac{2}{3} = \frac{1}{3}.
  \]
  
  在 \(\triangle ABC\) 中，\(A+B+C = \pi \implies C = \pi - (A+B) = \pi - S\)。
  
  根据诱导公式：
  \[
    \cos C = \cos(\pi - S) = -\cos S = -\frac{1}{3}.
  \]
  故选 A。%
}

% =====================================================
% 别名兼容：旧课次及学生个人宏定义
% =====================================================
\newcommand{\qTJWTriangleSumToProductStem}{\qTriangleSumToProductStem}
\newcommand{\qTJWTriangleSumToProductOptions}{\qTriangleSumToProductOptions}
\newcommand{\qTJWTriangleSumToProductAnswer}{\qTriangleSumToProductAnswer}
\newcommand{\qTJWTriangleSumToProductSolution}{\qTriangleSumToProductSolution}
